\documentclass[11pt,aspectratio=169,notes,mathserif]{beamer}
\usepackage[utf8]{inputenc}
\usepackage[T1]{fontenc}
\usefonttheme{serif}
\setbeamertemplate{footnote}{\href{https://EnergyAdvisor.app}{\beamergeftableotenote{~}}}
\setbeamertemplate{footnote}{\justifying}
\usepackage[english]{babel}
\usetheme{Madrid}
\usepackage{graphicx}
\usepackage{bm}
\usepackage{ragged2e}
\setbeamertemplate{caption}[numbered]
\usepackage{amsmath}
\usepackage{amsfonts}
\usepackage{amssymb}
\usepackage{lmodern}
\usepackage[T1]{fontenc}
\usepackage[utf8]{inputenc}
\usepackage{hyperref}
\usepackage{listings}
\usepackage{colortbl}

% Logo in top right corner
\addtobeamertemplate{frametitle}{}{%
    \begin{tikzpicture}[remember picture,overlay]
        \node[anchor=north east, xshift=-0.6cm, yshift=-0.2cm] at (current page.north east) {
            \includegraphics[width=0.8cm]{college_logo}
        };
    \end{tikzpicture}
}

% For images
\usepackage{graphicx}
\graphicspath{{./images/}}

% Justification
\setbeamerfont{footline}{size=\fontsize{8}{8}\selectfont}
\usepackage{ragged2e}
\apptocmd{\frame}{}{\justifying}{}

% Title page adjustments
\defbeamertemplate{title page}{custom}[1][]{
    \begin{centering}
        \vspace*{1cm}
        {\usebeamerfont{title}\usebeamercolor[fg]{title}\textbf{#1}\par}%
        \vskip1em
        {\usebeamerfont{subtitle}\usebeamercolor[fg]{subtitle}{#1}\par}%
        \vskip2em
        \begin{columns}
            \column{1\linewidth}
        \end{columns}
    \end{centering}
}

\title[]{%
    \textbf{Energy Efficiency Advisor} \\[0.5cm]
    \textit{AI-Powered Energy Consumption Analysis \& Recommendations}
}

\author[]{%
    \textbf{Team Members:} \\
    \vspace{0.3cm}
    {\small Your Name 1 \\ Roll No: XXXXXXXX} \\[0.2cm]
    {\small Your Name 2 \\ Roll No: XXXXXXXX} \\[0.2cm]
    {\small Your Name 3 \\ Roll No: XXXXXXXX} \\[0.2cm]
    {\small Your Name 4 \\ Roll No: XXXXXXXX}
}

\institute[]{%
    \textbf{Under the Guidance of} \\[0.3cm]
    \textbf{Dr. [Guide Name]} \\
    \textbf{Associate Professor,} \\
    \textbf{Dept of [Your Department],} \\
    \textbf{[Your College Name]}
}

\date{\today}

\begin{document}
% Title Frame
\begin{frame}[plain]
    \maketitle
\end{frame}

% Table of Contents
\begin{frame}{\textbf{Table of Content}}
    \tableofcontents
\end{frame}

\section{Abstract}
\begin{frame}{\textbf{Abstract}}
    \justifying
    Energy Efficiency Advisor is an AI-powered web application designed to help Indian households analyze their energy consumption patterns. 
    
    The application allows users to input their household appliances and usage patterns, then uses Google Gemini AI to generate personalized recommendations for reducing energy consumption.
    
    \begin{block}{\textbf{Key Features}}
        \begin{itemize}
            \item \textbf{18 Common Appliances} - Covers all major household devices
            \item \textbf{Usage-Based Analysis} - Light/Normal/Heavy usage patterns
            \item \textbf{AI-Powered Recommendations} - Personalized tips from Gemini AI
            \item \textbf{No Technical Knowledge Required} - Simple, intuitive interface
        \end{itemize}
    \end{block}
    
    This project demonstrates the practical application of AI in addressing real-world energy consumption challenges faced by everyday households.
\end{frame}

\section{Problem Statement}
\begin{frame}{\textbf{Problem Statement}}
    \justifying
    Indian households face several challenges related to energy consumption:
    
    \begin{block}{\textbf{Challenges}}
        \begin{itemize}
            \item \textbf{Rising Electricity Costs} - Average household spends ₹3,000-8,000/month
            \item \textbf{Lack of Awareness} - Users don't know which appliances consume the most power
            \item \textbf{No Personalized Advice} - Generic tips don't account for individual usage
            \item \textbf{Technical Complexity} - Existing solutions require technical expertise
            \item \textbf{Difficulty Identifying Waste} - Hard to pinpoint energy-wasting habits
        \end{itemize}
    \end{block}
    
    \begin{block}{\textbf{Need}}
        A simple, AI-powered tool that provides personalized, actionable recommendations for reducing energy consumption without requiring technical knowledge.
    \end{block}
\end{frame}

\section{Introduction}
\begin{frame}{\textbf{Introduction}}
    \justifying
    \textbf{Energy Efficiency Advisor} is a modern solution to household energy management.
    
    \begin{block}{\textbf{What It Does}}
        \begin{itemize}
            \item Analyzes household energy consumption based on appliances and usage
            \item Uses AI to identify energy waste
            \item Provides personalized, actionable recommendations
            \item Shows potential savings in kWh (not money)
        \end{itemize}
    \end{block}
    
    \begin{block}{\textbf{Why India-Focused?}}
        \begin{itemize}
            \item Tailored for Indian household appliances
            \item Common usage patterns in Indian homes
            \item Focus on kWh savings rather than currency
            \item Addresses specific challenges of Indian climate and lifestyle
        \end{itemize}
    \end{block}
\end{frame}

\begin{frame}{\textbf{Introduction (Cont.)}}
    \justifying
    The application transforms complex energy calculations into simple, actionable insights.
    
    \begin{columns}
        \column{0.5\linewidth}
        \begin{block}{\textbf{Before}}
            \begin{itemize}
                \item High electricity bills
                \item No idea what's wasting power
                \item Generic, useless advice
                \item Complex technical terms
            \end{itemize}
        \end{block}
        
        \column{0.5\linewidth}
        \begin{block}{\textbf{After}}
            \begin{itemize}
                \item Clear usage breakdown
                \item Specific waste identified
                \item Personalized tips
                \item Simple action steps
            \end{itemize}
        \end{block}
    \end{columns}
    
    \begin{figure}
        \centering
        \includegraphics[width=0.7\linewidth]{app_interface}
        \caption{Energy Efficiency Advisor Interface}
    \end{figure}
\end{frame}

\section{Objectives}
\begin{frame}{\textbf{Objectives}}
    \justifying
    The primary objectives of this project are:
    
    \begin{block}{\textbf{Primary Objectives}}
        \begin{enumerate}
            \item \textbf{Simplify Energy Analysis} - Make complex energy calculations accessible to everyone without technical knowledge
            
            \item \textbf{Identify Energy Waste} - Highlight which appliances and usage patterns waste the most energy
            
            \item \textbf{Provide Actionable Recommendations} - Give step-by-step advice categorized by difficulty (Easy/Medium/Hard)
            
            \item \textbf{Reduce Energy Consumption} - Help users save kWh with practical, implementable changes
            
            \item \textbf{India-Focused Solution} - Tailor the solution for Indian households, appliances, and usage patterns
        \end{enumerate}
    \end{block}
\end{frame}

\section{Technology Stack}
\begin{frame}{\textbf{Technology Stack}}
    \justifying
    
    \begin{table}
        \centering
        \small
        \caption{Technologies Used}
        \begin{tabular}{|c|p{5cm}|p{4cm}|}
            \hline
            \textbf{Component} & \textbf{Technology} & \textbf{Purpose} \\ \hline
            Frontend & React 18 + Vite & Modern UI framework \\ \hline
            Styling & CSS & Custom styling \\ \hline
            AI Model & Google Gemini 2.5 Flash & AI recommendations \\ \hline
            API & REST (Direct) & Avoids CORS issues \\ \hline
            Deployment & Static Hosting & Easy deployment \\ \hline
        \end{tabular}
    \end{table}
    
    \begin{block}{\textbf{Why These Technologies?}}
        \begin{itemize}
            \item \textbf{React + Vite} - Fast, component-based, hot module replacement
            \item \textbf{Gemini 2.5 Flash} - Fast AI processing, affordable, high quality
            \item \textbf{Direct REST API} - No SDK dependencies, avoids CORS
        \end{itemize}
    \end{block}
\end{frame}

\begin{frame}{\textbf{Architecture Overview}}
    \justifying
    
    \begin{figure}
        \centering
        \includegraphics[width=0.85\linewidth]{architecture}
        \caption{System Architecture}
    \end{figure}
    
    \begin{block}{\textbf{Data Flow}}
        User Input $\rightarrow$ React Form $\rightarrow$ Gemini API $\rightarrow$ AI Analysis $\rightarrow$ Recommendations
    \end{block}
\end{frame}

\section{Implementation}
\subsection{Modules Used}
\begin{frame}{\textbf{Implementation: Modules Used}}
    \justifying
    
    \begin{block}{\textbf{1. Appliance Database (appliances.js)}}
        \begin{itemize}
            \item 18 pre-configured household appliances
            \item Default wattage values based on real data
            \item Usage presets (Light/Normal/Heavy hours)
            \item Icons for visual representation
        \end{itemize}
    \end{block}
    
    \begin{block}{\textbf{2. Home Form Component (HomeForm.jsx)}}
        \begin{itemize}
            \item Household info inputs (occupants, location)
            \item Interactive appliance selection grid
            \item Usage configuration per appliance
            \item Additional context textarea
        \end{itemize}
    \end{block}
    
    \begin{block}{\textbf{3. Gemini Service (gemini.js)}}
        \begin{itemize}
            \item Prompt construction from user data
            \item Direct API integration with Gemini
            \item JSON response parsing
            \item Error handling
        \end{itemize}
    \end{block}
\end{frame}

\subsection{Working Process}
\begin{frame}{\textbf{Implementation: Working Process}}
    \justifying
    
    \begin{block}{\textbf{Step 1: User Input}}
        \begin{itemize}
            \item Enter household info (occupants, location)
            \item Select applicable appliances from 18 options
            \item Configure usage level (Light/Normal/Heavy)
            \item Add optional context (work from home, etc.)
        \end{itemize}
    \end{block}
    
    \begin{block}{\textbf{Step 2: Data Processing}}
        \begin{itemize}
            \item Calculate monthly kWh for each appliance
            \item Construct AI prompt with all data
            \item Format data for Gemini API
        \end{itemize}
    \end{block}
    
    \begin{block}{\textbf{Step 3: AI Analysis}}
        \begin{itemize}
            \item Send request to Gemini 2.5 Flash API
            \item AI analyzes consumption patterns
            \item Generate waste analysis
            \item Create saving recommendations
        \end{itemize}
    \end{block}
\end{frame}

\begin{frame}{\textbf{Implementation: Working Process (Cont.)}}
    \justifying
    
    \begin{block}{\textbf{Step 4: Results Display}}
        \begin{itemize}
            \item Show total monthly usage in kWh
            \item Highlight biggest energy consumers
            \item Display identified waste analysis
            \item List actionable saving tips with difficulty levels
        \end{itemize}
    \end{block}
    
    \begin{figure}
        \centering
        \includegraphics[width=0.75\linewidth]{results_screenshot}
        \caption{Results Display}
    \end{figure}
\end{frame}

\begin{frame}{\textbf{Implementation: Sample Code}}
    \justifying
    
    \begin{block}{\textbf{Gemini API Integration}}
        \small
        \begin{verbatim}
const API_URL = 'https://generativelanguage.googleapis.com
                  /v1beta/models/gemini-2.5-flash:generateContent';

const response = await fetch(`${API_URL}?key=${API_KEY}`, {
  method: 'POST',
  headers: { 'Content-Type': 'application/json' },
  body: JSON.stringify({
    contents: [{ parts: [{ text: prompt }] }]
  })
});
        \end{verbatim}
    \end{block}
    
    \begin{block}{\textbf{Appliance Data Structure}}
        \small
        \begin{verbatim}
{
  id: 'ac',
  name: 'Air Conditioning',
  icon: '❄️',
  wattage: 3500,
  usageOptions: [
    { label: 'Off', hours: 0 },
    { label: 'Few hours (4h)', hours: 4 },
    { label: 'Normal (8h)', hours: 8 },
  ]
}
        \end{verbatim}
    \end{block}
\end{frame}

\section{Results}
\begin{frame}{\textbf{Results}}
    \justifying
    
    \begin{block}{\textbf{Analysis Output}}
        \begin{itemize}
            \item \textbf{Total Monthly Usage} - Calculated in kWh
            \item \textbf{Biggest Consumer} - Identifies top energy user
            \item \textbf{Waste Analysis} - Shows where energy is being wasted
            \item \textbf{Saving Tips} - Actionable advice per appliance
            \item \textbf{Reduction Estimate} - Percentage savings possible
        \end{itemize}
    \end{block}
    
    \begin{figure}
        \centering
        \includegraphics[width=0.65\linewidth]{results_summary}
        \caption{Results Summary View}
    \end{figure}
\end{frame}

\begin{frame}{\textbf{Results (Cont.)}}
    \justifying
    
    \begin{block}{\textbf{Sample Recommendations Generated}}
        \begin{enumerate}
            \item \textbf{AC Usage} - "Turn off AC when not in room" - Save 120 kWh/month
            \item \textbf{Lighting} - "Switch to LED bulbs" - Save 15 kWh/month
            \item \textbf{Chargers} - "Unplug when not in use" - Save 8 kWh/month
            \item \textbf{Refrigerator} - "Check door seals" - Save 20 kWh/month
        \end{enumerate}
    \end{block}
    
    \begin{block}{\textbf{Metrics Achieved}}
        \begin{itemize}
            \item \textbf{Total Appliances Analyzed} - 18 common devices
            \item \textbf{Potential Reduction} - 15-25\% of current usage
            \item \textbf{Response Time} - Few seconds per analysis
            \item \textbf{User Interface} - Clean, intuitive design
        \end{itemize}
    \end{block}
\end{frame}

\section{Conclusion}
\begin{frame}{\textbf{Conclusion}}
    \justifying
    
    \begin{block}{\textbf{Summary}}
        \begin{itemize}
            \item Successfully built an AI-powered energy advisor
            \item Simplified complex energy analysis for common users
            \item Focused on Indian household context and appliances
            \item Provided actionable, personalized recommendations
            \item Demonstrated practical application of AI in daily life
        \end{itemize}
    \end{block}
    
    \begin{block}{\textbf{Key Achievements}}
        \begin{enumerate}
            \item Clean, intuitive React interface
            \item Integration with Google Gemini AI
            \item 18 appliances with realistic wattage values
            \item Complete analysis pipeline from input to results
            \item No technical knowledge required to use
        \end{enumerate}
    \end{block}
\end{frame}

\begin{frame}{\textbf{Future Scope}}
    \justifying
    
    \begin{block}{\textbf{Potential Enhancements}}
        \begin{itemize}
            \item \textbf{Save History} - Store analysis in localStorage
            \item \textbf{Export Reports} - Generate PDF reports
            \item \textbf{Scenario Comparison} - Compare multiple setups
            \item \textbf{Progress Tracking} - Monitor savings over time
            \item \textbf{More Appliances} - Expand appliance database
            \item \textbf{Location-Based} - Weather and climate insights
            \item \textbf{Mobile App} - Native iOS/Android version
        \end{itemize}
    \end{block}
    
    \begin{block}{\textbf{Impact}}
        This project demonstrates how AI can be used to solve everyday problems and help ordinary people make better decisions about energy consumption.
    \end{block}
\end{frame}

\section{References}
\begin{frame}{\textbf{References}}
    \begin{thebibliography}{3}
        \setbeamertemplate{bibliography item}[article]
        \justifying
        
        \bibitem{googleAI}
        Google AI, \textit{Gemini API Documentation}, Google, 2024. 
        [Online]. Available: \href{https://ai.google.dev/docs}{ai.google.dev/docs}
        
        \bibitem{react}
        React Team, \textit{React Documentation}, Meta Platforms, 2024.
        [Online]. Available: \href{https://react.dev}{react.dev}
        
        \bibitem{vite}
        Vite Team, \textit{Vite Documentation}, 2024.
        [Online]. Available: \href{https://vitejs.dev}{vitejs.dev}
        
    \end{thebibliography}
\end{frame}

\begin{frame}{\textbf{Thank You}}
    \begin{center}
        \Large
        \textbf{Thank You!}
        
        \vspace{1cm}
        \Large
        \textbf{Questions?}
        
        \vspace{2cm}
        \normalsize
        \textit{Energy Efficiency Advisor} \\
        \small AI-Powered Energy Consumption Analysis
        
    \end{center}
\end{frame}

\end{document}
